\documentclass[../../../include/open-logic-section]{subfiles}

\begin{document}

\olfileid[tr]{sth}{spine}{valpha}

\olsection{Aşamaları $V_\alpha$'lar Olarak Tanımlamak}

\olref[sth][ordinals][]{chap} içinde iyi sıralamaları ve (von Neumann)
sıra sayılarını tanımladık. Bu bölümde bunları küme hiyerarşisinin
\emph{kendisini} nitelemek için kullanacağız. Bunun için
\olref[sth][ordinals][opps]{sec} içinde ardıl ve limit sıra sayıları
fikirlerini tanımladığımızı hatırlayın. Bu fikirleri aşağıdaki tanımda
kullanıyoruz:

\begin{defn}\ollabel{defValphas}
	\begin{align*}
	V_\emptyset &\defis \emptyset\\
	V_{\ordsucc{\alpha}} &\defis \Pow{V_\alpha} & &
	\text{her $\alpha$ sıra sayısı için}\\
	V_{\alpha} &\defis \bigcup_{\gamma < \alpha} V_\gamma & &
	\text{$\alpha$ bir limit sıra sayısı olduğunda.}
\end{align*}
\end{defn}
\noindent
Bu, sıra sayıları üzerinde \emph{sonlu ötesi özyinelemeyle} yapılmış bir
tanım olacaktır. Bu bakımdan doğal sayılar üzerindeki fonksiyonların
özyinelemeli tanımlarıyla karşılaştırılmalıdır.\footnote{Toplama, çarpma ve
üs alma tanımları için \olref[sfr][infinite][dedekind]{sec} kesimine
bakın.} Doğal sayılarda olduğu gibi bir başlangıç durumu ile ardıl durumlar
tanımlanır; fakat sıra sayılarıyla çalışırken \emph{limit} durumlarındaki
davranışı da betimlememiz gerekir.

$V_\alpha$'ların bu tanımı önemli bir dönüm noktası olacaktır. Küme
hiyerarşimizi, kümeleri \emph{aşamalar} hâlinde oluşturmakla gayriresmî
olarak gerekçelendirdik. $V_\alpha$'lar gerçekte tam da bu aşamalardır.
Ancak önemli olarak bu, aşamaların \emph{içsel} bir nitelenmesidir. Teorinin
olası bir \emph{modelini} önermek yerine aşamaları küme teorimizin
\emph{içinde} tanımlamış olacağız.

\end{document}
